\documentclass[12pt]{article}
\usepackage[utf8]{inputenc}
\usepackage[T1]{fontenc}
\usepackage{hyperref}
\title{Asmat, Casuarina Coast Language Resources for K-12 Curriculum Development}
\author{Compiled by Claude Sonnet 5.5 (claude-sonnet-5-5) for Andr\'as Kornai}
\date{}
\begin{document}
\maketitle
\section{Introduction}
Casuarina Coast Asmat (ISO 639-3 asc) is an endangered Trans-New Guinea language of Papua, Indonesia. It is not known to be taught in schools. Web searches found very little material beyond one scholarly description, wordlists, online lookup tools and scripture; no school materials were found.
\section{Primary Educational Resources}
No primers, textbooks or school programmes in the language were found in the searches.
\section{Linguistic and Reference Materials}
\begin{itemize}
\item Voorhoeve, C. L. 1980. \emph{The Asmat Languages of Irian Jaya} (Pacific Linguistics B-64); bibliographic record: \url{https://glottolog.org/resource/reference/id/85998}
\item Glosbe online entries (automatically assembled, not a vetted dictionary): \url{https://glosbe.com/asc/en}
\item TransNewGuinea.org comparative wordlists (example entry): \url{https://transnewguinea.org/word/crab}
\item Overview of the Asmat languages: \url{https://www.Wikipedia.com/wiki/Asmat_language}
\end{itemize}
\section{Religious and Community Resources}
\begin{itemize}
\item Scripture and Global Recordings listings: \url{https://scriptureearth.org/00eng.php?iso=asc}
\end{itemize}
\section{Digital and Audio Resources}
Audio recordings are listed via the ScriptureEarth page above. No other digital resources were found.
\section{Curriculum Development Context}
Any curriculum would have to be built largely from scratch, drawing on Voorhoeve's description and on community work with speakers. The sources found are old or unvetted, so new documentation is needed.
\end{document}